% PDF pp. 376–384
\chapter*{索引}
\addcontentsline{toc}{chapter}{索引}
\markboth{索引}{}

以下页码是英文原版的印刷页码。\textbf{粗体}页码指表格，\textit{斜体}页码指图。词项保留原索引的字母分组和层级。下表列出两版章首页，可点击中文章名跳到相应章节。

\begin{center}
\begin{tabular}{lrrr}
  \toprule
  中文章节 & 原 PDF 页序 & 原版印刷页 & 中文页\\
  \midrule
  \hyperref[ch:01]{第 1 章\quad 光滑流形} & 18  & 1   & \pageref*{ch:01}\\
  \hyperref[ch:02]{第 2 章\quad 黎曼几何} & 58  & 41  & \pageref*{ch:02}\\
  \hyperref[ch:03]{第 3 章\quad 微分形式} & 99  & 82  & \pageref*{ch:03}\\
  \hyperref[ch:04]{第 4 章\quad 李导数}   & 164 & 147 & \pageref*{ch:04}\\
  \hyperref[ch:05]{第 5 章\quad 黎曼几何进阶专题}
                                      & 223 & 206 & \pageref*{ch:05}\\
  \hyperref[ch:06]{第 6 章\quad 流形上的统计理论}
                                      & 272 & 255 & \pageref*{ch:06}\\
  \bottomrule
\end{tabular}
\end{center}

\newcommand{\IndexLetter}[1]{%
  \par\addvspace{.55\baselineskip}\needspace{2\baselineskip}%
  \noindent\textbf{#1}\par\nobreak\smallskip
}
\newcommand{\IndexEntry}[3]{%
  \par\begingroup
  \setlength{\leftskip}{#1em}%
  \noindent\hangindent=1em\hangafter=1
  #2\if\relax\detokenize{#3}\relax\else，#3\fi\par
  \endgroup
}
\setlength{\columnsep}{17bp}
\begin{multicols}{2}
  \small
  \raggedcolumns
  \raggedright
  % Index terms in the source order, with the original printed page locators.
% PDF p. 376
\IndexLetter{A}
\IndexEntry{0}{伴随}{\mbox{350--354}}
\IndexEntry{0}{仿射联络}{\mbox{300}}
\IndexEntry{1}{坐标表示}{\mbox{301--303}}
\IndexEntry{1}{定义}{\mbox{303}}
\IndexEntry{1}{方向导数}{\mbox{300--301}}
\IndexEntry{0}{曲面上的角与球面上的角}{\mbox{242}}
\IndexEntry{0}{反对称性}{\mbox{95}}
\IndexEntry{0}{阿司匹林}{\mbox{191}}
\IndexEntry{1}{图表示}{\mbox{182}，\mbox{\textbf{183}}}
\IndexEntry{1}{反应位点预测}{\mbox{190}}
\IndexEntry{0}{渐近行为}{\mbox{262}}
\IndexEntry{0}{渐近分布}{}
\IndexEntry{1}{与检验统计量}{\mbox{281}，\mbox{286--287}，\mbox{335--337}}
\IndexEntry{1}{用德尔塔法推导}{\mbox{336}，\mbox{338--339}}
\IndexEntry{0}{单元格比例向量的渐近正态性}{\mbox{335}}
\IndexEntry{0}{流形上统计量的渐近理论}{\mbox{255--256}}
\IndexEntry{0}{图册}{\mbox{3}}
\IndexLetter{B}
\IndexEntry{0}{后向条件期望}{\mbox{340--341}}
\IndexEntry{0}{柯尔莫哥洛夫后向方程}{\mbox{340}}
\IndexEntry{0}{柯尔莫哥洛夫后向偏微分方程}{\mbox{353--354}}
\IndexEntry{0}{流形上的二项分布}{\mbox{263--265}}
\IndexEntry{0}{生化系统中的因果网络推断}{\mbox{202}}
\IndexEntry{0}{生物系统}{\mbox{196}}
\IndexEntry{0}{二重向量}{\mbox{93}}
\IndexEntry{0}{带边界流形}{\mbox{11--12}}
\IndexEntry{0}{边界点}{\mbox{12}}
\IndexEntry{0}{边值问题}{\mbox{132--133}}
\IndexEntry{0}{布朗运动}{\mbox{288}}
\IndexLetter{C}
\IndexEntry{0}{黑色素瘤}{\mbox{202--203}}
\IndexEntry{0}{癌症研究}{}
\IndexEntry{1}{图神经网络}{\mbox{202}}
\IndexEntry{1}{反应–扩散模型}{\mbox{202}}
\IndexEntry{0}{嘉当魔术公式}{\mbox{124}，\mbox{126}}
\IndexEntry{0}{带曲率修正的单元格比例概率}{\mbox{333--334}}
\IndexEntry{0}{单元格比例向量的渐近正态性}{\mbox{335}}
\IndexEntry{0}{中心极限定理（CLT）}{}
\IndexEntry{1}{弗雷歇均值的中心极限定理}{\mbox{258--259}}
\IndexEntry{1}{黎曼中心极限定理}{\mbox{314}}
\IndexEntry{0}{CGR，见混沌游戏表示（CGR）}{}
\IndexEntry{0}{流形上的换元公式}{\mbox{276}}
\IndexEntry{0}{混沌游戏表示（CGR）}{\mbox{27}，\mbox{285}}
\IndexEntry{1}{受扩散模型启发的扩展}{\mbox{30--33}}
\IndexEntry{1}{纳入基因组位置信息}{\mbox{29--30}}
\IndexEntry{1}{迭代更新规则}{\mbox{29}}
\IndexEntry{1}{SNP 基因型的 CGR}{\mbox{29}}
\IndexEntry{0}{坐标卡引理}{\mbox{20--21}}
\IndexEntry{0}{坐标卡}{}
\IndexEntry{1}{基本概念}{\mbox{2--3}}
\IndexEntry{1}{计算方法}{\mbox{4--5}}
\IndexEntry{1}{连续函数的图像}{\mbox{3}，\mbox{\textit{3}}}
\IndexEntry{1}{球面}{\mbox{4}，\mbox{5}}
\IndexEntry{0}{陈–高斯–博内定理}{\mbox{228--229}}
\IndexEntry{0}{卡方分布}{}
\IndexEntry{1}{曲率和非线性}{\mbox{277--278}}
\IndexEntry{1}{流形上的卡方分布}{\mbox{277}}
\IndexEntry{0}{克里斯托费尔符号}{\mbox{246}，\mbox{299}}
\IndexEntry{1}{一阶导数}{\mbox{321--324}}
\IndexEntry{1}{列维-奇维塔联络的公式}{\mbox{307--308}}
\IndexEntry{0}{经典斯托克斯定理}{\mbox{108--109}，\mbox{113}}
\IndexEntry{0}{CLT，见中心极限定理（CLT）}{}
\IndexEntry{0}{计算过渡映射}{\mbox{7--8}}
\IndexEntry{0}{生成元与条件漂移}{\mbox{340--342}}
\IndexEntry{0}{守恒形式}{\mbox{352}}
\IndexEntry{0}{常曲率流形}{\mbox{330}}
\IndexEntry{0}{常正曲率}{\mbox{283}}
\IndexEntry{0}{列联表}{\mbox{279}，\mbox{\textbf{279}}，\mbox{281--282}，\mbox{\textbf{282}}}
\IndexEntry{1}{从基因型到 FCGR 和列联表}{\mbox{285}}
\IndexEntry{1}{统计量}{\mbox{285--287}}
\IndexEntry{0}{缩并}{\mbox{139}}
\IndexEntry{0}{逆变分量}{\mbox{317--318}}
\IndexEntry{0}{逆变指标}{\mbox{317}}
\IndexEntry{0}{坐标形式}{\mbox{352}}
\IndexEntry{0}{坐标限制公式}{\mbox{144--145}}
\IndexEntry{0}{充分统计量的协方差矩阵}{\mbox{270}}
\IndexEntry{0}{切空间中的协方差算子 $\Sigma$}{\mbox{313--315}}
% PDF p. 377
\IndexEntry{0}{协变分量}{\mbox{318}}
\IndexEntry{0}{协变导数}{\mbox{243}}
\IndexEntry{1}{沿曲线的协变导数}{\mbox{232}}
\IndexEntry{1}{克里斯托费尔符号公式（列维-奇维塔联络）}{\mbox{307--308}}
\IndexEntry{1}{交换性}{\mbox{245}，\mbox{304--305}}
\IndexEntry{2}{交换公式}{\mbox{242--243}}
\IndexEntry{1}{定义}{\mbox{305}}
\IndexEntry{1}{一般公式}{\mbox{305--307}，\mbox{309--311}}
\IndexEntry{1}{交换协变导数}{\mbox{234--236}}
\IndexEntry{1}{性质}{\mbox{305}}
\IndexEntry{0}{协变指标}{\mbox{317}}
\IndexEntry{0}{余向量场的李导数}{\mbox{157--158}，\mbox{164--166}，\mbox{170--172}}
\IndexEntry{0}{双曲消息传递网络的分类交叉熵损失}{\mbox{78}}
\IndexEntry{0}{叉积}{}
\IndexEntry{1}{应用}{\mbox{91--92}}
\IndexEntry{1}{几何解释}{\mbox{91}}
\IndexEntry{1}{数学定义}{\mbox{90--91}}
\IndexEntry{1}{性质}{\mbox{91}}
\IndexEntry{1}{用楔积和霍奇星算子表示}{\mbox{92--95}}
\IndexEntry{0}{曲率}{\mbox{80}}
\IndexEntry{1}{单元格比例概率的曲率修正}{\mbox{333--334}}
\IndexEntry{1}{可积性}{\mbox{225}}
\IndexEntry{0}{曲率张量}{\mbox{244--245}}
\IndexEntry{1}{坐标公式}{\mbox{249}}
\IndexEntry{1}{球面的曲率张量}{\mbox{247--248}}
\IndexLetter{D}
\IndexEntry{0}{DDPM，见去噪扩散概率模型（DDPM）}{}
\IndexEntry{0}{变形理论}{\mbox{163}}
\IndexEntry{0}{用德尔塔法推导渐近分布}{\mbox{336}，\mbox{338--339}}
\IndexEntry{0}{去噪扩散概率模型（DDPM）}{\mbox{32}}
\IndexEntry{0}{流形上指数族的密度函数}{\mbox{268--270}}
\IndexEntry{0}{确定性常微分方程采样}{\mbox{298}}
\IndexEntry{0}{微分同胚}{\mbox{42}}
\IndexEntry{0}{反向时间随机微分方程}{\mbox{339--342}}
\IndexEntry{0}{微分形式}{\mbox{82--83}}
\IndexEntry{1}{积分的数学公式}{\mbox{100--108}}
\IndexEntry{0}{微分 $1$-形式（直观介绍）}{\mbox{84--86}}
\IndexEntry{1}{与力场关联的 $1$-形式}{\mbox{85}}
\IndexEntry{1}{计算功}{\mbox{85--86}}
\IndexEntry{1}{微分函数}{\mbox{86}}
\IndexEntry{1}{1-形式的直观解释}{\mbox{85}}
\IndexEntry{1}{直观解释}{\mbox{87--88}}
\IndexEntry{1}{拉回}{\mbox{150--151}}
\IndexEntry{1}{光滑曲线}{\mbox{86}}
\IndexEntry{0}{微分几何}{\mbox{8}}
\IndexEntry{1}{高斯映射}{\mbox{214--217}}
\IndexEntry{1}{过渡映射}{\mbox{5--6}}
\IndexEntry{0}{$\R^n$ 上的微分 $k$-形式}{\mbox{95--100}}
\IndexEntry{1}{定义}{\mbox{96}}
\IndexEntry{1}{数学表达式}{\mbox{96}}
\IndexEntry{1}{切空间}{\mbox{96}}
\IndexEntry{0}{比较不同点处的向量}{\mbox{304}}
\IndexEntry{0}{菲克扩散定律}{\mbox{199--200}}
\IndexEntry{0}{受扩散模型启发的 CGR 扩展}{\mbox{30--33}}
\IndexEntry{0}{扩散映射}{}
\IndexEntry{1}{用于降维}{\mbox{295--296}}
\IndexEntry{1}{用于提升}{\mbox{296--297}}
\IndexEntry{0}{基于得分的扩散模型}{\mbox{291}}
\IndexEntry{0}{黎曼流形上的扩散–反应方程}{\mbox{200--201}}
\IndexEntry{0}{扩散情形中的伴随 $L_t^*$}{\mbox{350--354}}
\IndexEntry{0}{用于降维的扩散映射}{\mbox{295--296}}
\IndexEntry{0}{离散算子与双曲图拉普拉斯算子}{\mbox{136}}
\IndexEntry{0}{流形上随机向量变换的分布理论}{\mbox{272--276}}
\IndexEntry{0}{流形上的散度定理}{\mbox{127}}
\IndexEntry{1}{欧氏空间中的散度}{\mbox{128}}
\IndexEntry{1}{流形上的向量场}{\mbox{127}}
\IndexEntry{0}{双重扩散映射}{\mbox{296--297}}
\IndexEntry{0}{药物发现}{\mbox{191--195}}
\IndexEntry{0}{对偶余标架}{}
\IndexEntry{1}{标架变换下的对偶余标架}{\mbox{45--46}}
\IndexEntry{1}{与标架}{\mbox{45}}
\IndexLetter{E}
\IndexEntry{0}{边级李导数}{\mbox{187}，\mbox{189}}
\IndexEntry{0}{用特征分解求加权卡方分布}{\mbox{336--337}}
\IndexEntry{0}{爱因斯坦求和约定}{\mbox{44--45}}
\IndexEntry{0}{流形上基于能量距离的双样本检验}{\mbox{261--262}}
\IndexEntry{0}{能量距离统计量}{\mbox{261}}
\IndexEntry{0}{欧氏情形}{\mbox{272}}
\IndexEntry{1}{渐近分布}{\mbox{279--280}}
\IndexEntry{0}{欧氏几何}{\mbox{54}}
\IndexEntry{0}{欧氏度量与极坐标}{\mbox{48--49}}
\IndexEntry{0}{欧氏空间}{\mbox{247}，\mbox{263}}
\IndexEntry{1}{散度}{\mbox{128}}
\IndexEntry{1}{图拉普拉斯算子}{\mbox{135--136}}
\IndexEntry{1}{$\R^n$}{\mbox{14}}
\IndexEntry{0}{指数族}{}
% PDF p. 378
\IndexEntry{1}{表述}{\mbox{271--272}}
\IndexEntry{1}{流形上的指数族}{\mbox{268--270}}
\IndexEntry{0}{指数映射}{}
\IndexEntry{1}{定义}{\mbox{70}}
\IndexEntry{1}{一般黎曼流形中的公式}{\mbox{70--71}}
\IndexEntry{1}{双曲面（洛伦兹）模型中}{\mbox{66--70}}
\IndexEntry{1}{与雅可比场}{\mbox{324}}
\IndexEntry{1}{庞加莱球模型中的指数映射}{\mbox{73--74}}
\IndexEntry{1}{投影回双曲空间}{\mbox{75}}
\IndexEntry{1}{由指数映射拉回的度量}{\mbox{324--325}}
\IndexEntry{1}{与黎曼正规坐标}{\mbox{324}}
\IndexEntry{1}{从切空间到流形}{\mbox{276}}
\IndexEntry{0}{推广的高斯–博内公式}{\mbox{226--228}}
\IndexEntry{0}{外微分}{}
\IndexEntry{1}{定义}{\mbox{98}}
\IndexEntry{1}{几何解释}{\mbox{99--100}}
\IndexEntry{1}{性质}{\mbox{98--99}}
\IndexLetter{F}
\IndexEntry{0}{FCGR 矩阵，见频率混沌游戏表示（FCGR）矩阵}{}
\IndexEntry{0}{菲克扩散定律}{\mbox{199--200}}
\IndexEntry{0}{有限维向量空间}{\mbox{12--13}}
\IndexEntry{0}{首次返回映射，见庞加莱映射}{}
\IndexEntry{0}{福克–普朗克方程（柯尔莫哥洛夫前向方程）}{\mbox{339}，\mbox{340}}
\IndexEntry{0}{与力场关联的微分 $1$-形式（直观介绍）}{\mbox{85}}
\IndexEntry{0}{前向扩散（受 CGR 启发的噪声）}{\mbox{31}}
\IndexEntry{0}{前向生成元}{\mbox{340}}
\IndexEntry{0}{标架}{}
\IndexEntry{1}{与对偶余标架}{\mbox{45}}
\IndexEntry{1}{标架变换下的黎曼度量}{\mbox{48--49}}
\IndexEntry{0}{弗雷歇函数}{\mbox{257}}
\IndexEntry{0}{弗雷歇均值}{}
\IndexEntry{1}{中心极限定理}{\mbox{258--259}}
\IndexEntry{1}{定义}{\mbox{257}}
\IndexEntry{1}{直观解释和例子}{\mbox{257}}
\IndexEntry{0}{频率混沌游戏表示（FCGR）矩阵}{\mbox{284--288}}
\IndexEntry{1}{用格拉斯曼流形分析}{\mbox{33--36}}
\IndexEntry{1}{构造}{\mbox{28}}
\IndexEntry{0}{弗罗贝尼乌斯定理}{\mbox{163}}
\IndexEntry{0}{满秩矩阵}{\mbox{16}}
\IndexLetter{G}
\IndexEntry{0}{高斯–博内定理}{\mbox{222--223}}
\IndexEntry{0}{黎曼流形上的高斯–格林（散度）定理}{\mbox{131}}
\IndexEntry{0}{高斯近似}{\mbox{277}}
\IndexEntry{0}{高斯曲率}{}
\IndexEntry{1}{常用计算公式}{\mbox{220--222}}
\IndexEntry{1}{定义}{\mbox{220}}
\IndexEntry{0}{各向同性高斯密度}{\mbox{273}}
\IndexEntry{0}{高斯映射}{}
\IndexEntry{1}{微分}{\mbox{215--216}}
\IndexEntry{1}{在微分几何中}{\mbox{214--217}}
\IndexEntry{1}{测地三角形的高斯映射}{\mbox{239--242}}
\IndexEntry{1}{测地线在高斯映射下的像}{\mbox{242}}
\IndexEntry{1}{形状算子}{\mbox{215--216}}
\IndexEntry{1}{单位法向量}{\mbox{215}}
\IndexEntry{0}{基因表达建模中的李导数}{\mbox{198--203}}
\IndexEntry{0}{一般指数映射}{\mbox{73}}
\IndexEntry{0}{广义格拉斯曼距离}{\mbox{34--35}}
\IndexEntry{0}{一般线性群}{\mbox{15--16}}
\IndexEntry{0}{广义斯托克斯定理}{\mbox{111--113}}
\IndexEntry{0}{一般张量场的李导数}{\mbox{168--177}}
\IndexEntry{0}{生成模型（GM）}{\mbox{290}}
\IndexEntry{0}{生成元}{}
\IndexEntry{1}{与条件漂移}{\mbox{340--342}}
\IndexEntry{1}{定义}{\mbox{347--349}}
\IndexEntry{1}{伊藤扩散的生成元}{\mbox{346--350}}
\IndexEntry{0}{作为动力系统的基因调控（常微分方程形式）}{\mbox{201}}
\IndexEntry{0}{基因调控网络（GRNs）}{\mbox{196--197}}
\IndexEntry{0}{一般流形}{\mbox{71}}
\IndexEntry{0}{基因型}{}
\IndexEntry{1}{SNP 基因型的混沌游戏表示}{\mbox{29}}
\IndexEntry{1}{从基因型到 FCGR 和列联表}{\mbox{285}}
\IndexEntry{0}{测地线}{}
\IndexEntry{1}{定义}{\mbox{232}}
\IndexEntry{1}{测地偏离}{\mbox{244}}
\IndexEntry{1}{测地距离}{\mbox{58--63}，\mbox{261}，\mbox{337--338}}
\IndexEntry{1}{测地线方程}{\mbox{230--234}}
\IndexEntry{1}{测地线模型}{\mbox{58}}
\IndexEntry{1}{性质}{\mbox{231--232}}
\IndexEntry{1}{变分}{\mbox{229--230}}
\IndexEntry{0}{测地偏离定律}{\mbox{230}}
\IndexEntry{0}{测地距离（格拉斯曼距离）}{\mbox{33--34}，\mbox{337--339}}
\IndexEntry{0}{双曲拉普拉斯算子的测地极坐标形式}{\mbox{135}}
\IndexEntry{0}{几何代数}{\mbox{95}}
\IndexEntry{0}{几何控制理论}{\mbox{163--164}}
\IndexEntry{0}{几何感知亲和权重}{\mbox{136}}
\IndexEntry{0}{几何感知解码器}{\mbox{298}}
\IndexEntry{0}{几何感知图神经网络}{\mbox{177}}
\IndexEntry{0}{吉拉尔定理}{\mbox{241--242}}
\IndexEntry{0}{胰高血糖素受体（GCGR）抑制剂}{\mbox{192}}
\IndexEntry{0}{GM，见生成模型（GM）}{}
\IndexEntry{0}{GNNs，见图神经网络（GNNs）}{}
\IndexEntry{0}{欧氏空间中的图拉普拉斯算子}{\mbox{135--136}}
% PDF p. 379
\IndexEntry{0}{图神经网络（GNNs）}{\mbox{74}，\mbox{180}，\mbox{189}}
\IndexEntry{1}{在癌症研究中}{\mbox{202}}
\IndexEntry{1}{几何感知图神经网络}{\mbox{177}}
\IndexEntry{1}{李导数增强的图神经网络}{\mbox{182}}
\IndexEntry{1}{与机器学习}{\mbox{177}}
\IndexEntry{1}{消息传递}{\mbox{194}}
\IndexEntry{1}{连续流建模}{\mbox{195--196}}
\IndexEntry{1}{创造性与新颖性建模}{\mbox{193--194}}
\IndexEntry{1}{随时间和空间变化的图}{\mbox{193}}
\IndexEntry{0}{图神经反应–扩散模型}{\mbox{202}}
\IndexEntry{0}{思维图}{\mbox{147}}
\IndexEntry{0}{通过正交投影描述格拉斯曼几何}{\mbox{24}}
\IndexEntry{0}{格拉斯曼流形}{}
\IndexEntry{1}{基于格拉斯曼流形的方法}{\mbox{286--287}}
\IndexEntry{1}{格拉斯曼距离}{\mbox{33}}
\IndexEntry{1}{用于分析 FCGR}{\mbox{33--36}}
\IndexEntry{1}{在格拉斯曼框架下考察数据}{\mbox{285}}
\IndexEntry{1}{格拉斯曼流形上的测地距离}{\mbox{337--339}}
\IndexEntry{1}{作为商空间}{\mbox{24}}
\IndexEntry{0}{格拉斯曼流形}{\mbox{22}，\mbox{24}，\mbox{39}}
\IndexEntry{1}{线性子空间的格拉斯曼流形}{\mbox{23}}
\IndexEntry{1}{表示}{\mbox{37--38}}
\IndexEntry{0}{GRNs，见基因调控网络（GRNs）}{}
\IndexLetter{H}
\IndexEntry{0}{豪斯多夫性}{\mbox{2}}
\IndexEntry{0}{热方程}{\mbox{133}}
\IndexEntry{0}{高维积分}{\mbox{106--108}}
\IndexEntry{0}{HMPNs，见双曲消息传递网络（HMPNs）}{}
\IndexEntry{0}{HNJ，见双曲邻接法（HNJ）}{}
\IndexEntry{0}{霍奇星算子}{\mbox{93--94}}
\IndexEntry{0}{同伦公式}{\mbox{126}}
\IndexEntry{0}{霍特林 $T^2$ 检验}{\mbox{259}}
\IndexEntry{0}{混合传播}{\mbox{186}}
\IndexEntry{0}{经双曲适配的多层感知机}{\mbox{75}}
\IndexEntry{0}{双曲圆盘}{}
\IndexEntry{1}{庞加莱模型}{\mbox{54}}
\IndexEntry{1}{里奇张量}{\mbox{252}}
\IndexEntry{0}{节点的双曲嵌入}{\mbox{74--75}，\mbox{136}}
\IndexEntry{0}{双曲几何}{}
\IndexEntry{1}{适配}{\mbox{75}}
\IndexEntry{1}{庞加莱圆盘模型}{\mbox{58--66}}
\IndexEntry{0}{离散算子与双曲图拉普拉斯算子}{\mbox{136}}
\IndexEntry{0}{双曲拉普拉斯算子的测地极坐标形式}{\mbox{135}}
\IndexEntry{0}{双曲拉普拉斯 Transformer（HL-Former）}{\mbox{136--139}}
\IndexEntry{0}{双曲流形}{\mbox{51}}
\IndexEntry{0}{双曲消息传递网络（HMPNs）}{\mbox{74}}
\IndexEntry{1}{聚合}{\mbox{76}}
\IndexEntry{1}{组合损失}{\mbox{78}}
\IndexEntry{1}{组成部分}{\mbox{74}}
\IndexEntry{1}{分类交叉熵损失}{\mbox{78}}
\IndexEntry{1}{双曲池化}{\mbox{77}}
\IndexEntry{1}{损失函数}{\mbox{77--78}}
\IndexEntry{1}{消息传递}{\mbox{76}}
\IndexEntry{1}{节点更新}{\mbox{77}}
\IndexEntry{0}{双曲中点}{\mbox{66}}
\IndexEntry{0}{双曲邻接法（HNJ）}{\mbox{65--66}}
\IndexEntry{0}{切向量的双曲范数}{\mbox{68}}
\IndexEntry{0}{双曲系统发育树的推断}{\mbox{63--66}}
\IndexEntry{0}{双曲平面}{\mbox{300}}
\IndexEntry{0}{双曲池化}{\mbox{77}}
\IndexEntry{0}{双曲位置编码（H-LPE）}{\mbox{137}}
\IndexEntry{0}{双曲空间}{}
\IndexEntry{1}{拉普拉斯–贝尔特拉米算子}{\mbox{133--135}}
\IndexEntry{1}{庞加莱球模型}{\mbox{78}}
\IndexEntry{1}{切空间}{\mbox{67--70}，\mbox{\textit{68}}，\mbox{\textit{69}}}
\IndexEntry{0}{双曲面}{}
\IndexEntry{1}{作为水平集}{\mbox{19}}
\IndexEntry{1}{洛伦兹模型}{\mbox{51--53}，\mbox{66--71}}
\IndexEntry{1}{类型}{\mbox{19}}
\IndexEntry{0}{超球面坐标卡（经纬度的推广）}{\mbox{79}}
\IndexEntry{0}{超曲面上的诱导曲面测度 $dS_g$}{\mbox{144}}
\IndexEntry{0}{假设检验}{\mbox{263}}
\IndexEntry{1}{与统计推断}{\mbox{267}}
\IndexEntry{1}{流形上的检验}{\mbox{255}}
\IndexLetter{I}
\IndexEntry{0}{等距浸入}{\mbox{49}}
\IndexEntry{0}{示性函数}{\mbox{264}}
\IndexEntry{0}{诱导黎曼度量}{\mbox{80}}
\IndexEntry{0}{信息几何}{\mbox{255}}
\IndexEntry{0}{胰高血糖素受体抑制剂}{\mbox{192}}
\IndexEntry{0}{沿曲线积分 $1$-形式}{\mbox{101--102}}
\IndexEntry{0}{在曲面上积分 $2$-形式}{\mbox{102--106}}
\IndexEntry{0}{在 $k$ 维流形上积分 $k$-形式}{\mbox{106--108}}
\IndexEntry{0}{微分形式积分的数学公式}{\mbox{100--108}}
\IndexEntry{0}{线性随机微分方程的积分因子法}{\mbox{354--356}}
\IndexEntry{0}{内边处的抵消}{\mbox{225}}
\IndexEntry{0}{内点}{\mbox{11}}
\IndexEntry{0}{内积算子}{\mbox{139--141}}
\IndexEntry{0}{等距}{\mbox{44}}
\IndexEntry{0}{等距浸入}{\mbox{49}}
\IndexEntry{0}{各向同性高斯密度}{\mbox{273}}
% PDF p. 380
\IndexEntry{0}{各向同性潜空间噪声}{\mbox{298}}
\IndexEntry{0}{伊藤积分}{}
\IndexEntry{1}{背景与直观解释}{\mbox{343}}
\IndexEntry{1}{定义}{\mbox{343--344}}
\IndexEntry{1}{直接使用等距性}{\mbox{345}}
\IndexEntry{1}{基本性质}{\mbox{344--345}}
\IndexEntry{1}{乘积法则}{\mbox{346}}
\IndexLetter{J}
\IndexEntry{0}{雅可比矩阵}{\mbox{56}，\mbox{122}}
\IndexEntry{0}{雅可比方程}{\mbox{229}，\mbox{230}，\mbox{233--236}}
\IndexEntry{1}{求解算法}{\mbox{236--238}}
\IndexEntry{1}{应用}{\mbox{238--239}}
\IndexEntry{1}{用于求解的截面曲率}{\mbox{325--326}}
\IndexEntry{1}{用于求解的泰勒展开}{\mbox{326--328}}
\IndexEntry{0}{雅可比场}{\mbox{229}}
\IndexEntry{1}{定义}{\mbox{229--230}}
\IndexEntry{1}{与指数映射}{\mbox{324}}
\IndexEntry{1}{用雅可比场表示二阶导数}{\mbox{325}}
\IndexLetter{K}
\IndexEntry{0}{$k$-形式}{\mbox{93}}
\IndexEntry{0}{柯尔莫哥洛夫后向方程}{\mbox{340}}
\IndexEntry{0}{柯尔莫哥洛夫前向方程／福克–普朗克方程}{\mbox{351--352}}
\IndexEntry{0}{$k$-向量}{\mbox{93}}
\IndexLetter{L}
\IndexEntry{0}{拉普拉斯–贝尔特拉米算子}{}
\IndexEntry{1}{双曲空间上}{\mbox{133--135}}
\IndexEntry{1}{庞加莱球中}{\mbox{135}}
\IndexEntry{1}{黎曼流形上}{\mbox{132--133}}
\IndexEntry{0}{拉普拉斯方程}{\mbox{132}}
\IndexEntry{0}{拉普拉斯扩散注意力偏置}{\mbox{137}}
\IndexEntry{0}{大语言模型（LLMs）}{\mbox{74}}
\IndexEntry{0}{大型推理模型（LRMs）}{\mbox{74}}
\IndexEntry{0}{经纬度的推广}{\mbox{79}}
\IndexEntry{0}{学习参数模型}{\mbox{32}}
\IndexEntry{0}{列维-奇维塔联络}{\mbox{235}}
\IndexEntry{0}{李括号}{\mbox{243}，\mbox{315--316}}
\IndexEntry{0}{李导数}{\mbox{156--157}}
\IndexEntry{1}{自适应加权}{\mbox{186}}
\IndexEntry{1}{优点}{\mbox{180}}
\IndexEntry{1}{应用}{\mbox{177--203}}
\IndexEntry{2}{自主系统}{\mbox{179}}
\IndexEntry{2}{生物学}{\mbox{178--179}}
\IndexEntry{2}{计算机视觉}{\mbox{178}}
\IndexEntry{2}{数据科学与数据分析}{\mbox{179}}
\IndexEntry{2}{微分几何}{\mbox{178}}
\IndexEntry{2}{经济学与博弈论}{\mbox{179}}
\IndexEntry{2}{工程应用}{\mbox{178}}
\IndexEntry{2}{几何与物理}{\mbox{177}}
\IndexEntry{2}{机器学习与图神经网络}{\mbox{177}}
\IndexEntry{2}{预测准确性}{\mbox{187--188}}
\IndexEntry{2}{机器人与控制系统}{\mbox{178}}
\IndexEntry{1}{计算公式}{\mbox{181}}
\IndexEntry{1}{余向量场的李导数}{\mbox{157--158}，\mbox{170--172}}
\IndexEntry{2}{$1$-形式}{\mbox{164--166}}
\IndexEntry{1}{定义}{\mbox{157}}
\IndexEntry{1}{不同类型张量场的李导数}{\mbox{160--167}}
\IndexEntry{1}{方向导数}{\mbox{157}}
\IndexEntry{1}{在动力学中的应用}{\mbox{196}}
\IndexEntry{1}{基于李导数的节点重要性框架}{\mbox{188}}
\IndexEntry{1}{在基因表达建模中的应用}{\mbox{198--203}}
\IndexEntry{1}{一般张量场的李导数}{\mbox{168--177}}
\IndexEntry{2}{一般性质}{\mbox{168--169}}
\IndexEntry{2}{数学表述}{\mbox{168}}
\IndexEntry{1}{在信息流中的应用}{\mbox{188}}
\IndexEntry{1}{积分}{\mbox{183--185}}
\IndexEntry{1}{归一化李导数}{\mbox{186}}
\IndexEntry{1}{普通导数}{\mbox{157}}
\IndexEntry{1}{可能的改法}{\mbox{186--187}}
\IndexEntry{1}{性质}{\mbox{159--160}}
\IndexEntry{1}{正则化李导数}{\mbox{186}，\mbox{189}，\mbox{191}}
\IndexEntry{1}{作用}{\mbox{166--167}}
\IndexEntry{1}{标量函数的李导数}{\mbox{160--161}}
\IndexEntry{1}{张量场的李导数}{\mbox{158--159}，\mbox{172--177}}
\IndexEntry{1}{向量场的李导数}{\mbox{161--162}，\mbox{169--170}，\mbox{201--202}}
\IndexEntry{0}{李导数增强的图神经网络（GNN）}{\mbox{182}}
\IndexEntry{0}{线性代数}{\mbox{65}}
\IndexEntry{0}{线性映射}{\mbox{42}}
\IndexEntry{0}{推导线性随机微分方程的解}{\mbox{354--356}}
\IndexEntry{0}{线性子空间的格拉斯曼流形}{\mbox{23}}
\IndexEntry{0}{曲线积分}{\mbox{101--102}}
\IndexEntry{0}{LLMs，见大语言模型（LLMs）}{}
\IndexEntry{0}{局部注意力机制}{\mbox{187}}
\IndexEntry{0}{局部坐标基}{\mbox{48}}
\IndexEntry{0}{局部对称性与几何}{\mbox{162--163}}
\IndexEntry{0}{对数映射}{}
\IndexEntry{1}{庞加莱球模型中的对数映射}{\mbox{73--74}}
\IndexEntry{1}{到切空间的投影}{\mbox{75}}
\IndexEntry{0}{对数映射（指数映射的逆）}{\mbox{71--73}}
\IndexEntry{0}{对数配分函数}{\mbox{272}}
\IndexEntry{0}{洛伦兹内积}{\mbox{51}}
\IndexEntry{0}{洛伦兹模型}{\mbox{51--53}，\mbox{71}}
\IndexEntry{0}{具有低维内在结构的密度的采样}{\mbox{295--296}}
\IndexEntry{0}{下指标（协变指标）}{\mbox{317}，\mbox{319--320}}
\IndexEntry{0}{LRMs，见大型推理模型（LRMs）}{}
% PDF p. 381
\IndexLetter{M}
\IndexEntry{0}{机器学习}{}
\IndexEntry{1}{算法}{\mbox{263}}
\IndexEntry{1}{与图神经网络}{\mbox{177}}
\IndexEntry{0}{流形}{\mbox{1}}
\IndexEntry{1}{流形上分布函数的应用}{\mbox{263}}
\IndexEntry{1}{流形上检验统计量的渐近分布}{\mbox{281}}
\IndexEntry{1}{流形上统计量的渐近理论}{\mbox{255--256}}
\IndexEntry{1}{流形上的二项分布}{\mbox{263--265}}
\IndexEntry{1}{带边界流形}{\mbox{11--12}}
\IndexEntry{1}{流形上的换元公式}{\mbox{276}}
\IndexEntry{1}{流形上的卡方分布}{\mbox{277--278}}
\IndexEntry{1}{流形上概率分布函数的推导}{\mbox{262--263}}
\IndexEntry{1}{流形上随机向量变换的分布理论}{\mbox{272--276}}
\IndexEntry{1}{流形上的经验分布函数}{\mbox{257--258}}
\IndexEntry{1}{流形上基于能量距离的双样本检验}{\mbox{261--262}}
\IndexEntry{1}{流形上的指数族}{\mbox{268--270}}
\IndexEntry{1}{从切空间到流形的指数映射}{\mbox{276}}
\IndexEntry{1}{推广到流形}{\mbox{272--273}}
\IndexEntry{1}{流形上的一般定义}{\mbox{264}}
\IndexEntry{1}{流形上的假设检验}{\mbox{255}}
\IndexEntry{1}{流形上分布函数的必要性}{\mbox{263}}
\IndexEntry{1}{流形上的负二项分布}{\mbox{265--267}}
\IndexEntry{1}{流形上的非光滑函数}{\mbox{10}}
\IndexEntry{1}{不构成拓扑流形的空间}{\mbox{2}}
\IndexEntry{1}{流形的划分}{\mbox{278}}
\IndexEntry{1}{流形上的泊松分布}{\mbox{267--268}}
\IndexEntry{1}{流形上的概率分布}{\mbox{262--284}}
\IndexEntry{1}{乘积流形}{\mbox{4}}
\IndexEntry{1}{流形上的样本统计量}{\mbox{256--257}}
\IndexEntry{1}{流形上的采样方法}{\mbox{256}}
\IndexEntry{1}{流形上的统计理论}{\mbox{255}}
\IndexEntry{1}{流形结构的光滑性检验}{\mbox{22}}
\IndexEntry{1}{拓扑流形}{\mbox{1}}
\IndexEntry{1}{流形上的变换理论}{\mbox{273--276}}
\IndexEntry{1}{流形上的双向列联表}{\mbox{278--284}}
\IndexEntry{1}{流形上经验分布的方差}{\mbox{258}}
\IndexEntry{0}{流形扩散变分自编码器（MD-VAE）}{\mbox{288--289}，\mbox{\textit{289}}}
\IndexEntry{1}{潜在密度建模与采样}{\mbox{290--291}}
\IndexEntry{1}{模型架构}{\mbox{289--290}}
\IndexEntry{1}{基于得分的扩散模型}{\mbox{291}}
\IndexEntry{1}{七类漂移和扩散}{\mbox{292--293}}
\IndexEntry{1}{训练目标}{\mbox{290}}
\IndexEntry{0}{流形假设}{\mbox{255}}
\IndexEntry{0}{流形上的设定}{\mbox{265}}
\IndexEntry{1}{从二项分布到负二项分布}{\mbox{266--267}}
\IndexEntry{1}{内在测度与区域}{\mbox{265}}
\IndexEntry{1}{随机变量}{\mbox{266}}
\IndexEntry{0}{流形加权 $X^2$ 与 $p$ 值}{\mbox{284}}
\IndexEntry{0}{满秩矩阵}{\mbox{16}}
\IndexEntry{0}{矩阵指数法}{\mbox{238}}
\IndexEntry{0}{MCS 估计，见基于蒙特卡罗模拟（MCS）的估计}{}
\IndexEntry{0}{MD-VAE，见流形扩散变分自编码器（MD-VAE）}{}
\IndexEntry{0}{充分统计量的均值}{\mbox{269--270}}
\IndexEntry{0}{均方误差（MSE）}{\mbox{294}}
\IndexEntry{0}{黑色素瘤}{\mbox{202--203}}
\IndexEntry{0}{消息传递}{\mbox{185}}
\IndexEntry{0}{消息传递神经网络（MPNN）}{\mbox{76}}
\IndexEntry{1}{其中的张量}{\mbox{180--181}}
\IndexEntry{1}{其中的向量场与张量}{\mbox{180--181}}
\IndexEntry{0}{度量相容性}{\mbox{307}}
\IndexEntry{0}{度量关系}{\mbox{318--320}}
\IndexEntry{0}{度量张量展开}{\mbox{329}}
\IndexEntry{0}{度量张量正则化}{\mbox{288}}
\IndexEntry{0}{MLP，见多层感知机（MLP）}{}
\IndexEntry{0}{基于蒙特卡罗模拟（MCS）的估计}{\mbox{293--295}}
\IndexEntry{0}{MPNN，见消息传递神经网络（MPNN）}{}
\IndexEntry{0}{多层感知机（MLP）}{\mbox{74}，\mbox{75}}
\IndexEntry{0}{多线性函数}{\mbox{95}}
\IndexEntry{0}{多位点遗传关联研究}{\mbox{284--288}}
\IndexEntry{0}{多项分布协方差}{}
\IndexEntry{1}{含曲率}{\mbox{333--335}}
\IndexEntry{1}{无曲率}{\mbox{334}}
\IndexLetter{N}
\IndexEntry{0}{$n$ 维拓扑流形}{\mbox{1}}
\IndexEntry{0}{流形上的负二项分布}{\mbox{265--267}}
\IndexEntry{0}{$\R^n$ 上的 $n$-形式}{}
\IndexEntry{1}{定义}{\mbox{88}}
\IndexEntry{1}{解释}{\mbox{89}}
\IndexEntry{1}{直观理解}{\mbox{89--90}}
\IndexEntry{1}{性质}{\mbox{89}}
\IndexEntry{0}{邻接法（NJ）}{\mbox{65}}
\IndexEntry{0}{流的不可交换性}{\mbox{162}}
\IndexEntry{0}{流形上的非光滑函数}{\mbox{10}}
\IndexEntry{0}{非拓扑流形}{\mbox{2}}
\IndexEntry{0}{黎曼正态分布}{\mbox{270--272}}
\IndexEntry{0}{归一化李导数}{\mbox{186}}
\IndexLetter{O}
\IndexEntry{0}{ODE，见常微分方程（ODE）}{}
\IndexEntry{0}{开子流形}{\mbox{14--15}}
% PDF p. 382
\IndexEntry{0}{常微分方程（ODE）}{\mbox{237--238}，\mbox{294}}
\IndexEntry{0}{正交}{\mbox{42}}
\IndexEntry{0}{正交投影算子}{\mbox{37}}
\IndexEntry{0}{用正交投影算子表示格拉斯曼几何}{\mbox{24}}
\IndexEntry{0}{标准正交}{\mbox{42}}
\IndexEntry{0}{标准正交切基下的矩阵表示}{\mbox{219}}
\IndexLetter{P}
\IndexEntry{0}{平行移动与旋转}{\mbox{244}}
\IndexEntry{0}{参数估计}{\mbox{263}}
\IndexEntry{0}{PDF，见概率分布函数（PDF）}{}
\IndexEntry{0}{皮尔逊卡方统计量}{\mbox{336}}
\IndexEntry{0}{作为二次型的皮尔逊 $X^2$}{\mbox{336}}
\IndexEntry{0}{置换程序}{\mbox{262}}
\IndexEntry{0}{物理人工智能}{\mbox{262}}
\IndexEntry{0}{依据物理定律的梯度}{\mbox{185}}
\IndexEntry{0}{庞加莱球模型}{}
\IndexEntry{1}{指数映射和对数映射}{\mbox{73--74}}
\IndexEntry{1}{双曲空间的球模型}{\mbox{78}}
\IndexEntry{1}{拉普拉斯–贝尔特拉米算子}{\mbox{135}}
\IndexEntry{1}{度量}{\mbox{310--311}}
\IndexEntry{0}{庞加莱圆盘模型}{\mbox{53--58}，\mbox{134}}
\IndexEntry{0}{庞加莱双曲圆盘}{\mbox{54}，\mbox{58--66}}
\IndexEntry{0}{庞加莱引理}{\mbox{114}，\mbox{142--144}}
\IndexEntry{1}{闭微分形式}{\mbox{115--118}}
\IndexEntry{1}{恰当微分形式}{\mbox{118--120}}
\IndexEntry{1}{同伦公式}{\mbox{121}}
\IndexEntry{1}{证明思路}{\mbox{121}}
\IndexEntry{1}{拉回}{\mbox{121--122}}
\IndexEntry{1}{引理的表述}{\mbox{120}}
\IndexEntry{0}{庞加莱映射}{\mbox{114}}
\IndexEntry{0}{流形上的泊松分布}{\mbox{267--268}}
\IndexEntry{0}{主曲率}{}
\IndexEntry{1}{计算算法}{\mbox{217--219}}
\IndexEntry{1}{作为特征值的主曲率}{\mbox{219--220}}
\IndexEntry{1}{法曲率}{\mbox{217}}
\IndexEntry{1}{主曲率}{\mbox{217}}
\IndexEntry{0}{流形上的概率分布函数（PDF）}{\mbox{262--263}}
\IndexEntry{0}{流形上的概率分布}{\mbox{262--284}}
\IndexEntry{0}{欧氏情形下的概率质量函数（PMF）}{\mbox{265}}
\IndexEntry{0}{乘积 $i_X\omega$}{\mbox{139--140}}
\IndexEntry{0}{乘积流形}{\mbox{4}}
\IndexEntry{0}{乘积黎曼流形}{\mbox{\textit{50}}，\mbox{50--51}}
\IndexEntry{0}{投影映射}{\mbox{17}}
\IndexEntry{0}{射影空间}{\mbox{4}}
\IndexEntry{0}{伪向量}{\mbox{93}}
\IndexEntry{0}{拉回}{\mbox{147--148}}
\IndexEntry{1}{定义}{\mbox{\textit{148}}，\mbox{148--150}}
\IndexEntry{1}{微分 1-形式的拉回}{\mbox{150--151}}
\IndexEntry{1}{一般计算算法}{\mbox{151}}
\IndexEntry{1}{点处的拉回}{\mbox{150}}
\IndexEntry{1}{张量场 $T$ 的拉回}{\mbox{153--156}}
\IndexEntry{0}{拉回度量}{}
\IndexEntry{1}{计算}{\mbox{327--328}}
\IndexEntry{1}{经指数映射的拉回}{\mbox{324--325}}
\IndexEntry{0}{推前测度}{\mbox{272}}
\IndexEntry{0}{推前向量}{\mbox{149}}
\IndexLetter{Q}
\IndexEntry{0}{二次型形式的皮尔逊 $X^2$}{\mbox{336}}
\IndexEntry{0}{商空间}{\mbox{1}，\mbox{2}，\mbox{25--27}}
\IndexEntry{1}{作为商空间的格拉斯曼流形}{\mbox{24}}
\IndexEntry{1}{直观理解}{\mbox{22--23}}
\IndexEntry{0}{商结构}{\mbox{36}}
\IndexLetter{R}
\IndexEntry{0}{流形上的随机向量}{\mbox{272--276}}
\IndexEntry{0}{反应–扩散模型}{\mbox{198--200}，\mbox{202}}
\IndexEntry{0}{反应–扩散系统}{\mbox{200}}
\IndexEntry{0}{实双曲空间}{\mbox{51--62}}
\IndexEntry{0}{正则化李导数}{\mbox{189}，\mbox{191}}
\IndexEntry{0}{正则值定理}{\mbox{17--18}，\mbox{20}}
\IndexEntry{0}{强化学习（RL）中的双曲拉普拉斯 Transformer}{\mbox{138--139}}
\IndexEntry{0}{反向（去噪）扩散}{\mbox{31--32}}
\IndexEntry{0}{反向时间坐标}{\mbox{339--340}}
\IndexEntry{0}{反向时间随机微分方程}{\mbox{339--342}}
\IndexEntry{0}{里奇张量}{\mbox{250--251}}
\IndexEntry{1}{双曲圆盘的里奇张量}{\mbox{252}}
\IndexEntry{0}{黎曼曲率张量}{\mbox{243--244}，\mbox{317--318}，\mbox{320}}
\IndexEntry{0}{黎曼 Adam}{\mbox{66}}
\IndexEntry{0}{黎曼中心极限定理}{\mbox{314}}
\IndexEntry{0}{黎曼几何}{\mbox{329}}
\IndexEntry{1}{基本形式与主曲率}{\mbox{209--214}}
\IndexEntry{2}{第一基本形式}{\mbox{210}}
\IndexEntry{2}{第二基本形式}{\mbox{210--213}}
\IndexEntry{1}{黎曼几何基本定理}{\mbox{207--209}}
\IndexEntry{1}{预备知识与记号}{\mbox{207--208}}
\IndexEntry{0}{黎曼梯度}{\mbox{65}}
\IndexEntry{0}{黎曼对数映射}{}
\IndexEntry{1}{定义}{\mbox{311}}
\IndexEntry{1}{直观解释}{\mbox{312--313}}
% PDF p. 383
\IndexEntry{0}{黎曼流形}{\mbox{44}，\mbox{128}，\mbox{262}}
\IndexEntry{1}{流形上体积的定义}{\mbox{328}}
\IndexEntry{1}{黎曼流形上的扩散–反应方程}{\mbox{200--201}}
\IndexEntry{1}{一般黎曼流形的指数映射公式}{\mbox{70--71}}
\IndexEntry{1}{黎曼流形上的高斯–格林（散度）定理}{\mbox{131}}
\IndexEntry{1}{黎曼流形上的拉普拉斯–贝尔特拉米算子}{\mbox{132--133}}
\IndexEntry{1}{乘积黎曼流形}{\mbox{\textit{50}}，\mbox{50--51}}
\IndexEntry{1}{用于黎曼流形数据的双样本检验}{\mbox{259--262}}
\IndexEntry{0}{黎曼度量}{\mbox{41--44}，\mbox{\textit{43}}，\mbox{47--48}}
\IndexEntry{1}{换标架}{\mbox{48--49}}
\IndexEntry{1}{坐标表示}{\mbox{47}}
\IndexEntry{1}{诱导的黎曼度量}{\mbox{80}}
\IndexEntry{1}{局部表示}{\mbox{44--45}}
\IndexEntry{1}{$n$ 维球面的黎曼度量}{\mbox{251--252}}
\IndexEntry{1}{张量记法}{\mbox{46--48}}
\IndexEntry{0}{黎曼正态密度}{\mbox{270--271}}
\IndexEntry{0}{黎曼正态分布}{\mbox{270--272}}
\IndexEntry{0}{黎曼子流形}{\mbox{49}，\mbox{206--207}}
\IndexEntry{0}{指数映射与黎曼正规坐标}{\mbox{324}}
\IndexEntry{0}{“肿瘤”区域中的 RNA 速率}{\mbox{164}}
\IndexLetter{S}
\IndexEntry{0}{标量函数的李导数}{\mbox{160--161}}
\IndexEntry{0}{基于得分的扩散模型（SGM）}{\mbox{289}}
\IndexEntry{0}{得分匹配估计}{\mbox{295}}
\IndexEntry{0}{SDE，见随机微分方程（SDE）}{}
\IndexEntry{0}{SGM，见基于得分的扩散模型（SGM）}{}
\IndexEntry{0}{形状算子矩阵}{\mbox{216--217}}
\IndexEntry{0}{带符号面积密度}{\mbox{100}}
\IndexEntry{0}{单核苷酸多态性（SNP）基因型}{\mbox{28--29}}
\IndexEntry{0}{奇异值分解（SVD）}{\mbox{33}}
\IndexEntry{0}{光滑函数}{\mbox{8--10}}
\IndexEntry{1}{函数图像}{\mbox{16--17}}
\IndexEntry{1}{流形上的光滑函数}{\mbox{9--10}}
\IndexEntry{0}{光滑流形}{\mbox{4}}
\IndexEntry{1}{坐标卡引理}{\mbox{20--21}}
\IndexEntry{1}{定义}{\mbox{5--6}}
\IndexEntry{1}{水平集}{\mbox{18--19}}
\IndexEntry{0}{光滑映射}{\mbox{10--11}}
\IndexEntry{0}{SNP，见单核苷酸多态性（SNP）}{}
\IndexEntry{0}{索菲斯·李}{\mbox{147}}
\IndexEntry{0}{矩阵空间}{\mbox{13}}
\IndexEntry{0}{谱滤波门}{\mbox{137}}
\IndexEntry{0}{球面}{\mbox{4}}
\IndexEntry{1}{坐标}{\mbox{79--80}}
\IndexEntry{1}{定义}{\mbox{79}}
\IndexEntry{0}{球坐标中的散度}{\mbox{129--131}}
\IndexEntry{0}{测地三角形的球面几何}{\mbox{239--242}}
\IndexEntry{0}{施蒂费尔流形}{\mbox{36}}
\IndexEntry{0}{随机微分方程（SDE）}{\mbox{294}，\mbox{297}，\mbox{347}}
\IndexEntry{1}{线性随机微分方程}{\mbox{354--356}}
\IndexEntry{1}{反向时间随机微分方程}{\mbox{339--342}}
\IndexEntry{0}{斯托克斯定理}{}
\IndexEntry{1}{经典形式}{\mbox{108--109}，\mbox{113}}
\IndexEntry{1}{广义形式}{\mbox{111--113}}
\IndexEntry{1}{证明}{\mbox{109--111}}
\IndexEntry{0}{浸没}{\mbox{17}}
\IndexEntry{0}{曲面积分}{\mbox{102--106}}
\IndexEntry{0}{SVD，见奇异值分解（SVD）}{}
\IndexLetter{T}
\IndexEntry{0}{切空间}{}
\IndexEntry{1}{切空间近似}{\mbox{277}}
\IndexEntry{1}{基}{\mbox{97}}
\IndexEntry{1}{切空间中的协方差算子 $\Sigma$}{\mbox{313--315}}
\IndexEntry{1}{定义}{\mbox{67}}
\IndexEntry{1}{双曲空间的切空间}{\mbox{67--70}，\mbox{\textit{68}}，\mbox{\textit{69}}}
\IndexEntry{1}{映到切空间}{\mbox{259--260}}
\IndexEntry{1}{切空间双样本检验}{\mbox{259--261}}
\IndexEntry{0}{切向量}{\mbox{149}}
\IndexEntry{1}{作为导子}{\mbox{97}}
\IndexEntry{1}{双曲范数}{\mbox{68}}
\IndexEntry{0}{求解雅可比方程的泰勒展开}{\mbox{326--328}}
\IndexEntry{0}{张量场}{}
\IndexEntry{1}{张量场的李导数}{\mbox{158--167}，\mbox{172--177}}
\IndexEntry{1}{张量场的拉回}{\mbox{153--156}，\mbox{\textit{154}}}
\IndexEntry{0}{张量性}{\mbox{208}}
\IndexEntry{0}{黎曼度量的张量记法}{\mbox{46--48}}
\IndexEntry{0}{张量积的定义}{\mbox{46--47}}
\IndexEntry{0}{消息传递神经网络中的张量}{\mbox{180--181}}
\IndexEntry{0}{曲率项}{\mbox{335}}
\IndexEntry{0}{检验统计量}{\mbox{338}}
\IndexEntry{0}{边界的拓扑不变性}{\mbox{12}}
\IndexEntry{0}{拓扑流形}{\mbox{1}，\mbox{15}}
\IndexEntry{0}{$X$ 上的拓扑}{\mbox{21--22}}
\IndexEntry{0}{作为水平集的环面}{\mbox{19}}
\IndexEntry{0}{简化数据验证}{\mbox{298}}
\IndexEntry{0}{五分类单元的简单例子}{\mbox{66}，\mbox{\textit{67}}}
\IndexEntry{0}{简单推理示例}{\mbox{137}，\mbox{139}}
\IndexEntry{0}{流形上的变换理论}{\mbox{273--276}}
\IndexEntry{0}{过渡映射}{\mbox{3}，\mbox{6--7}}
\IndexEntry{1}{计算}{\mbox{7--8}}
\IndexEntry{1}{微分几何中的过渡映射}{\mbox{5--6}}
\IndexEntry{1}{光滑性}{\mbox{22}}
\IndexEntry{0}{流形上的双向列联表}{\mbox{278--284}}
\IndexEntry{0}{第一类错误率}{\mbox{287--288}}
% PDF p. 384
\IndexLetter{U}
\IndexEntry{0}{上半空间模型}{\mbox{133--134}}
\IndexEntry{0}{上指标（逆变指标）}{\mbox{317}}
\IndexLetter{V}
\IndexEntry{0}{变分自编码器（VAE）}{\mbox{290}}
\IndexEntry{0}{向量场}{\mbox{315}}
\IndexEntry{1}{向量场的散度}{\mbox{128--131}}
\IndexEntry{1}{欧氏空间中的向量场}{\mbox{128}}
\IndexEntry{1}{向量场的李导数}{\mbox{161--162}，\mbox{169--170}，\mbox{201--202}}
\IndexEntry{1}{用向量场表示反应项}{\mbox{201}}
\IndexEntry{1}{消息传递神经网络中的向量场和张量}{\mbox{180--181}}
\IndexEntry{0}{向量与二重向量之间的映射}{\mbox{93--94}}
\IndexEntry{0}{有限维向量空间}{\mbox{12--13}}
\IndexEntry{0}{体积与行列式}{\mbox{90--95}}
\IndexEntry{0}{体积元展开}{\mbox{329--330}}
\IndexEntry{0}{体积展开}{\mbox{317--330}}
\IndexEntry{1}{体积展开与正规坐标}{\mbox{299--300}}
\IndexEntry{0}{小球的体积}{\mbox{330}}
\IndexEntry{1}{行列式展开}{\mbox{331}}
\IndexEntry{1}{度量展开}{\mbox{331}}
\IndexEntry{1}{平方根展开}{\mbox{331}}
\IndexEntry{1}{极坐标积分}{\mbox{331--333}}
\IndexEntry{1}{正规坐标下的设定}{\mbox{330--331}}
\IndexEntry{0}{冯·米塞斯–费希尔（vMF）模型}{\mbox{260}}
\IndexLetter{W}
\IndexEntry{0}{楔积（外积）}{}
\IndexEntry{1}{计算}{\mbox{94}}
\IndexEntry{1}{用楔积构造}{\mbox{97--98}}
\IndexEntry{1}{楔积与叉积}{\mbox{92--95}}
\IndexEntry{2}{二重向量与伪向量}{\mbox{93}}
\IndexEntry{2}{霍奇星算子}{\mbox{93--94}}
\IndexEntry{1}{定义}{\mbox{92}}
\IndexEntry{1}{几何解释}{\mbox{93}}
\IndexEntry{1}{性质}{\mbox{93}}

\end{multicols}
